\section{Deciding when a change of index preserves every finite colimit}
\label{sec:effective-finality}

Consider a calculation assembled from local finite sets, with specified maps between them. One would like to drop or duplicate some indexing objects without changing the result of gluing. A change of index can work for one particular diagram by accident. We give a uniformly primitive recursive test which answers the question for \emph{every} finite-set-valued diagram of the given shape. A negative answer supplies the diagram on which the comparison fails. The ordinary categorical criterion is the final-functor theorem (also called cofinality, Stacks Project, \href{https://stacks.math.columbia.edu/tag/09WN}{Tag 09WN}). Its uniform arithmetic content depends on how quotients, paths, and obstructions are represented.

The repository's \texttt{ipra\_categorical\_methods} develops finite diagram colimits by graph connectivity. The object-calculus module \texttt{Foundations/FiniteCoequalizer.lean} instead merges two labels per generating pair. Its header explicitly leaves identification of the fibres with the generated equivalence, compact indexing of classes, and the general diagram reduction untreated. We give a bounded proof for that very merge algorithm, then use the completed quotient in the index-change test.

\subsection{A quotient with paths and contiguous class numbers}

Write $[n]=\{0,\ldots,n-1\}$. Take a graph on $[N]$ with a list of $m$ oriented edges $(a_i,b_i)$. A zigzag is a vertex followed by a list of edge indices with a direction bit. The checker follows the indicated endpoints and verifies the final vertex. An edge may be traversed in either direction. Initialize $L_0(x)=x$. At step $i$, put $A=L_i(a_i)$ and $B=L_i(b_i)$ and replace each occurrence of $B$ by $A$ in the old label table. Denote the result by $L_{i+1}$. This is simultaneous replacement using the old $A,B$.

\begin{lemma}[Exact fibres of the sequential merge]
\label{lem:label-exact}
$L_m(x)=L_m(y)$ if and only if a zigzag joins $x$ and $y$. A PR function of the graph and the pair returns such a zigzag whenever the labels agree. Every label is a fixed point. The quotient can be numbered as $[r]$ with a PR section.
\end{lemma}
\begin{proof}
Keep a path $p_i(x)$ from $x$ to $L_i(x)$, initially empty. If $L_i(x)\ne B$, keep the path. Otherwise follow $p_i(x)$, then $p_i(b_i)$ in reverse, then the new edge from $b_i$ back to $a_i$, then $p_i(a_i)$. These endpoints match by the update condition. If all old paths have length at most $H_i$, the new bound is $H_{i+1}=3H_i+1$ with $H_0=0$. Thus $H_i=(3^i-1)/2$. This deliberately loose clock avoids a termination argument about shortest paths. The path tables and labels can be updated together by primitive recursion. For instance, lists of at most $2H_m$ signed edge indices are coded by length and fixed-width digits in base $2m+2$. Verification of all $N$ stored paths is a bounded test.

After step $i$ the endpoints of every edge already processed have equal labels. For the new edge both labels become $A$. For each previous edge apply the same replacement function to its two equal labels. The invariant is a bounded conjunction over the first $i$ edges. Consequently the final labels agree along any supplied zigzag, by induction over its successive vertices with its code fixed. Conversely, if the labels of $x$ and $y$ agree, concatenate $p_m(x)$ with the reverse of $p_m(y)$. This gives a path of length at most $2H_m$. Applying the forward implication to $p_m(x)$ gives $L_m(x)=L_m(L_m(x))$.

List in increasing order precisely the $v<N$ satisfying $L_m(v)=v$, say $s_0,\ldots,s_{r-1}$. Set $q(x)=\#\{v<L_m(x):L_m(v)=v\}$. The preceding fixed-point identity gives $s_{q(x)}=L_m(x)$, and counting the preceding list entries gives $q(s_k)=k$. A table $t:[N]\to[z]$ equal on the generating edges is constant along each supplied path; its unique factor is $k\mapsto t(s_k)$. The factor equation uses $p_m(x)$ and uniqueness follows by evaluating at $s_k$. If $N=0$ valid input forces $m=0$, and $r=0$ with empty structure maps. No positive-base digit operation is required in this case.
\end{proof}

A finite category $K$ has object and arrow lists, source and target tables, an identity table, and a composition table on composable pairs. The typing, identity and associativity checks scan at most triples of arrows. A finite-set diagram $D$ assigns $[n_d]$ to each object $d$ and a table $D(u):[n_d]\to[n_e]$ to each arrow $u:d\to e$, with unit and composition checked on entries. For its colimit take the $N=\sum_d n_d$ tagged elements $(d,x)$ in prefix-sum order. For each $u:d\to e$ and $x<n_d$ enter one edge
\[
 (d,x)\longrightarrow(e,D(u)x).
\]
There are $m=\sum_u n_{\mathrm{src}(u)}$ edges. A cocone to $[z]$ means tables $t_d:[n_d]\to[z]$ satisfying $t_e(D(u)x)=t_d(x)$. Concatenation identifies these exactly with the edge-invariant tables of the lemma. The compact quotient and the mediator just constructed therefore have the colimit universal property, including the empty diagram. Write $[d,x]_D$ for the class number of a tag.

\subsection{A complete test using representables}

For a finite functor $F:J\to K$ and $d\in K$, form a graph $G_d$. Its vertices are the pairs $(j,h)$ with $h:d\to F(j)$. For every $u:j\to j'$ add
\[
 (j,h)\longrightarrow (j',F(u)h).
\]
This is the underlying graph of the comma category $(d\mathbin\downarrow F)$. It has at most $|J_0|\,|K_1|$ vertices and $|J_1|\,|K_1|$ directed edges, where subscripts $0,1$ denote objects and arrows. Compute the number $r_d$ of components by Lemma~\ref{lem:label-exact}. Define $T(F)=1$ when every $r_d$ is one. For a diagram $D$ there is a comparison
\[
 c_D:\operatorname{colim}_{J}(DF)\longrightarrow\operatorname{colim}_{K}D,
 \qquad [j,x]_{DF}\longmapsto [F(j),x]_D .
\]
Each generating edge maps to a generating edge, proving that the formula respects classes. To implement it, use the section $s$ of the source quotient.

\begin{theorem}[Effective finite finality]
\label{thm:finite-finality}
For valid finite category and functor tables, iPRA proves that $T(F)=1$ is equivalent to bijectivity of $c_D$ for every valid numerical finite-set diagram code $D$ on $K$. On success the inverse is uniform PR in $(F,D)$. On failure a uniform PR procedure returns $d$ and a separating diagram $D=K(d,-)$, each of whose fibres has cardinality at most $|K_1|$. All quantified diagrams in this claim are numerical tables; their sizes vary without a preset bound.
\end{theorem}
\begin{proof}
In the successful case select the first vertex $v_d=(j_d,h_d)$ in each $G_d$. For $x\in D(d)$ send the tag $(d,x)$ to $[j_d,D(h_d)x]_{DF}$. For any other vertex $(j,h)$ the resulting class is the same: along an edge labelled $u$ the new value is $D(F(u)h)x=D(F(u))(D(h)x)$, and these two tags are connected by a generating edge of the restricted colimit. Reversing a step preserves equality. The path construction in Lemma~\ref{lem:label-exact} transports the equality to the selected vertex.

Now consider a generating edge $(d,x)\to(e,D(k)x)$ in the full diagram. In $G_d$ use the vertex $(j_e,h_e k)$. Its value agrees with the selected one and equals $[j_e,D(h_e)D(k)x]_{DF}$. Thus the proposed map is constant on the full diagram edges, so it descends to a map $b_D$ on its compact quotient. The composite $c_Db_D$ fixes $[d,x]_D$ by the defining edge for $h_d$. For $b_Dc_D$ on $[j,x]_{DF}$ take the vertex $(j,1_{F(j)})$ of $G_{F(j)}$ and use the vertex-independence already proved. Its value is $[j,x]_{DF}$ by the identity law. Both inverse equations hold pointwise on the complete quotient list.

If the bounded test fails, select its first offending $d$. Form $H_d(e)=K(d,e)$ by filtering the arrow list and identify it with the ordinal of its increasing ranks. An arrow $k:e\to e'$ acts by $h\mapsto kh$. The functor laws follow from identity and associativity. The graph for $\operatorname{colim}_K H_d$ is nonempty and every vertex $(e,h)$ is joined to $(d,1_d)$ by the edge for $h$; hence its quotient has one point. For the restricted diagram, the tagged vertices and edges are precisely those of $G_d$, up to the stated finite ranking. Its quotient has $r_d$ points. If $r_d=0$, the target's point has no preimage. If $r_d>1$, the first two distinct class numbers have the same image. The test distinguishes these cases by integer comparison. An assumed bijective comparison for every code would in particular apply to $H_d$, contradicting either witness. The initial split on the decidable value $T(F)$ gives the reverse implication intuitionistically.
\end{proof}

Here the apparently unbounded statement about all diagrams reduces to at most $|K_0|$ explicit test diagrams. Quantifier-free induction suffices: fix the arbitrary code $D$ during the forward proof. All identifications, edges, congruence steps along a path, and inverse equations are Boolean tests over its own finite tables. Reuse the displayed clocks, using parameters for the chosen tables and any path. The universal quantifier on $D$ is introduced after this proof. For $K$ empty a valid $F$ forces $J$ empty, and the empty comparison is a bijection.

\paragraph{A change invisible to a constant diagram.}
Take objects $s,t$, identities, two arrows $u,v:s\to t$, and an involution $g:t\to t$ interchanging $u,v$ by left composition. These are all the arrows. The inclusion of the full one-object subgroupoid at $t$ passes the test. Its comma graph over $s$ has vertices $u,v$ joined by the action of $g$; the graph over $t$ similarly joins $1_t,g$. Keeping only the identity at $t$ removes both connecting edges. The constant singleton diagram still has singleton colimit before and after restriction. But $K(s,-)$ has one element at $s$ and two at $t$, which become equal in the full colimit. The restriction now has two classes. The executable example checks all 18 diagrams on this $K$ whose object sizes are at most two: all comparisons pass for the full subgroupoid, while seven fail for the identity-only inclusion.

\paragraph{Transport to a presented target.}
The positive proof also extends a cocone for $DF$ in a target category $C$. For arrows $t_j:D(Fj)\to z$, set $\bar t_d=t_{j_d}D(h_d)$. A comma edge labelled $u$ preserves this expression by the cocone equation and associativity. Vertex-independence, proved along the supplied zigzags, establishes
$\bar t_d=\bar t_eD(k)$ by using $(j_e,h_e k)$ as above. Restriction recovers $t_j$ by the identity vertex. Extending the restriction of a full cocone $t$ recovers $t_d$ by its equation at $h_d$. Postcomposition by any $z\to z'$ commutes with this construction by associativity. Thus a universal cocone on the restriction extends to one on the full diagram, and conversely. This applies schematically for a fixed PR-presented target with proved laws and decidable hom equality. With arbitrary arithmetical hom equivalence one can also specialize to a fixed numerical $F$ and unroll its finitely many exhibited paths into equality steps. Uniform variable-length transport for a general nondecidable equivalence requires a separate derivation and is outside this conclusion.

\paragraph{Lean and the remaining object-calculus work.}
\texttt{categorica/Merge.lean} proves that a label agrees with another exactly for vertices joined by a generated zigzag, and proves idempotence. \texttt{categorica/Finality.lean} proves that a chosen comma vertex and paths extend a restricted set-valued cocone, with both pointwise inverse equations. Its \texttt{representable\_test} also proves that bijectivity on components for $K(d,-)$ is equivalent to connectedness and inhabitation of the comma graph, working entirely with representatives. These are complete semantic Lean lemmas. The third file, \texttt{PathDerivation.lean}, returns actual \(\Gamma\vdash_{\mathrm{iPRA}} t_x=t_y\) derivations from edge sequents when the external merge labels coincide. Its more general path-cut rule admits any formula-valued equivalence with supplied reflexivity, symmetry, and transitivity derivations. It includes a concrete reversed-edge instance. For a uniform internal derivation, use the existing \texttt{lab\_zero} and \texttt{lab\_succ} for $N>0$, handle $N=0$ separately, and code the oriented paths as above. Induct on the edge count with a single bounded bit testing that all $p_i(x)$ start and end correctly, use only indices below $i$, have length at most $H_i$, and all processed edges have equal labels. Then apply \texttt{FilteredEnumeration.sel\_rank} and \texttt{sel\_spec} to the fixed-label bit. The tag and comma tables are prefix sums, bounded filters, and finite composition lookups. Edge invariance is first proved for one entry, then along a supplied path by induction on its prefix length. Quotient factorization and the four displayed map equations finish the positive proof; filtering hom tables and comparing $r_d$ to one finish the converse. This plan yields an object theorem by threading the graph, table and diagram codes as numerical parameters of the PR terms. The coded uniform theorem remains a formalization target.
